% ============================================================
%  Chapter 3 -- Automata over Infinite Words
% ============================================================

\chapter{Automata over Infinite Words}
\label{ch:omega-automata}

\aside{Reactive systems do not produce
one finite output and stop.  Their
behaviour is an infinite interaction
with the environment.}

Finite-word automata are well suited for batch computations: a compiler
receives a finite program, a parser receives a finite string, a protocol
message has a finite body.  Verification of reactive systems is
different.  Operating systems, controllers, servers, and embedded
devices are meant to keep interacting with their environment.

Their behaviours are therefore not finite words but infinite words.  To
model them we keep the finite-state graph, but change the acceptance
condition.

\begin{chapteropening}
\textbf{Where we are.}  \Cref{ch:fsa,ch:regular-decision} developed the
finite-word theory: automata recognise finite behaviours, regular
expressions describe the same languages algebraically, and emptiness reduces
verification questions to graph reachability.

\textbf{What this chapter adds.}  We move from finite behaviours to infinite
executions.  The graph remains finite, but acceptance must now describe what
happens in the limit.

\textbf{Where it leads.}  B\"uchi automata become the automata-theoretic
model behind liveness, model checking, complementation, and temporal logic in
the following chapters.

\textbf{Chapter map.}
\begin{itemize}[nosep]
  \item \Cref{sec:omega-words} introduces infinite words as executions.
  \item \Cref{sec:omega-language-ops} lifts language operations to
        $\omega$-languages.
  \item \Cref{sec:buchi-automata} defines B\"uchi automata and their runs.
  \item \Cref{sec:omega-closure-expressions} gives closure constructions and
        $\omega$-regular expressions.
  \item \Cref{sec:buchi-emptiness} reduces non-emptiness to reachable lassos.
\end{itemize}
\end{chapteropening}

\constructionpipeline{2}

\begin{runningexample}{request/acknowledgement traces}
Throughout the infinite-word chapters, think of a reactive component that
observes requests and acknowledgements forever.  A typical liveness property is
``acknowledgement occurs infinitely often'', and the counterexamples are
ultimately periodic traces where acknowledgements eventually disappear.
\end{runningexample}


% ============================================================
\section{Infinite Words}
\label{sec:omega-words}
% ============================================================

Consider the request/acknowledgement component from the running example.
One possible execution is
\[
  \texttt{req}\ \texttt{ack}\ \texttt{req}\ \texttt{req}\ \texttt{ack}\cdots .
\]
The dots are not an omitted ending.  They are part of the model: the
component is expected to keep receiving inputs and producing observations
forever.  This is why we need words with positions $0,1,2,\ldots$ rather
than words with a last symbol.

\begin{definition}[$\omega$-word]
\label{def:omega-word}
Let $\Sigma$ be a finite alphabet.  An \textbf{$\omega$-word} over
$\Sigma$ is an infinite sequence
\[
  \alpha=\alpha_0\alpha_1\alpha_2\cdots
  \quad\text{with each }\alpha_i\in\Sigma.
\]
The set of all infinite words over $\Sigma$ is denoted by
$\Sigma^\omega$.
\end{definition}

An $\omega$-word is a function from the natural numbers to $\Sigma$:
position $0$ gets a symbol, position $1$ gets a symbol, and so on
without end.  The key operational difference from finite words is that
there is no ``last symbol'' to inspect, so every property of an
$\omega$-word must be phrased in terms of what happens eventually or
repeatedly, never in terms of how the word ends.

We use $u,v,w$ for finite words and Greek letters such as
$\alpha,\beta$ for infinite words.  This distinction matters because
finite words can be consumed completely, while infinite words have no
last position.

\begin{definition}[Segments and suffixes]
\label{def:omega-segments}
For an $\omega$-word $\alpha$, the finite segment
$\alpha[i,j]$ denotes the word
\[
  \alpha_i\alpha_{i+1}\cdots\alpha_{j-1}.
\]
The infinite suffix $\alpha[i,\omega]$ denotes
\[
  \alpha_i\alpha_{i+1}\alpha_{i+2}\cdots.
\]
\end{definition}

The convention $\alpha[i,i]=\varepsilon$ keeps formulas uniform: a
segment with equal endpoints consumes no symbols.

\begin{definition}[Infinitely often]
\label{def:inf-set}
For an infinite word $\alpha\in\Sigma^\omega$, define
\[
  \inf(\alpha)=\{a\in\Sigma : a\text{ occurs infinitely often in }\alpha\}.
\]
\end{definition}

Because $\Sigma$ is finite and $\alpha$ is infinite, at least one symbol
must occur infinitely often.  The set $\inf(\alpha)$ records the
long-run behaviour and forgets finite prefixes.

\begin{example}[A long-run property]
\label{ex:infinitely-many-a}
Let
\[
  \alpha = ab\,a\,bb\,a\,bbb\,a\,bbbb\cdots.
\]
Then $a$ occurs infinitely often and $b$ occurs infinitely often, so
$\inf(\alpha)=\{a,b\}$.
\end{example}


% ============================================================
\section{From Finite Languages to Infinite Languages}
\label{sec:omega-language-ops}
% ============================================================

We now have infinite words, but no way to describe sets of them.  Rather
than starting from scratch, we can lift familiar finite-word operations
to the infinite setting.  Some infinite languages can be built directly
from finite-word languages.

\begin{definition}[$\omega$-power]
\label{def:omega-power}
If $W\subseteq\Sigma^*$, then
\[
  W^\omega =
  \{w_0w_1w_2\cdots : w_i\in W\text{ for every }i\geq 0\}
\]
is the set of infinite concatenations of blocks from $W$.
\end{definition}

Usually we assume $\varepsilon\notin W$ when using $W^\omega$, because
empty blocks do not help construct an infinite word and can obscure the
intended decomposition.

\begin{definition}[Vectorial closure]
\label{def:vectorial-closure}
For $W\subseteq\Sigma^*$, define
\[
  \vec{W}=
  \{\alpha\in\Sigma^\omega :
    \alpha[0,n]\in W\text{ for infinitely many }n\}.
\]
\end{definition}

The language $\vec{W}$ contains the infinite words whose finite prefixes
hit $W$ infinitely often.

\begin{example}[Vectorial closure on the running example]
\label{ex:vectorial-closure-ack}
Over $\Sigma=\{\texttt{req},\texttt{ack}\}$, let $W$ be the set of
finite words that end with \texttt{ack}.  Then $\vec{W}$ is the set of
infinite traces such that infinitely many finite prefixes end with
\texttt{ack}---in other words, exactly the traces in which \texttt{ack}
occurs infinitely often.
\end{example}

The vectorial closure makes the connection to acceptance explicit.  It is a natural way to express recurrence:
not just ``a good prefix appears'', but ``good prefixes keep appearing''.
For a deterministic B\"uchi automaton this is exactly the operational
shape of acceptance: if $W$ is the set of finite prefixes that leave the
automaton in an accepting state, then the accepted $\omega$-language is
$\vec{W}$.  We will return to this observation in
\cref{ch:deterministic-automata-star-free}, where it becomes the precise
characterisation of deterministic B\"uchi languages.


% ============================================================
\section{B\"uchi Automata}
\label{sec:buchi-automata}
% ============================================================

A B\"uchi automaton has the same graph shape as an NFA---a finite set of
states, labelled transitions, an initial state, and a set of accepting
states---but it reads an infinite input rather than a finite one.  Since
the input never ends, there is no ``final state after the last symbol''.
Instead, the automaton asks a long-run question: does the run keep
returning to an accepting state, or does it eventually stay away
forever?  The tuple below captures exactly this setup: the graph is
finite, but the acceptance condition talks about what happens in the
limit.

\begin{definition}[B\"uchi automaton]
\label{def:buchi}
A \textbf{B\"uchi automaton} is a tuple
\[
  A=(Q,\Sigma,\Delta,q_0,F)
\]
where $Q$ is finite, $\Delta\subseteq Q\times\Sigma\times Q$,
$q_0\in Q$, and $F\subseteq Q$ is the set of accepting states.
\end{definition}

Read the tuple as a finite directed graph whose edges are labelled by
symbols: $Q$ is the set of nodes, $\Delta$ is the set of labelled edges,
$q_0$ marks the starting node, and $F$ marks the checkpoints that the
run must keep visiting.  Everything is identical to an NFA except the
role of $F$: in an NFA, $F$ says where a run must \emph{end}; in a
B\"uchi automaton, $F$ says where a run must \emph{return} infinitely
often.

\begin{definition}[Run and B\"uchi acceptance]
\label{def:buchi-run}
Let $\alpha=\alpha_0\alpha_1\alpha_2\cdots$.  A run of $A$ on
$\alpha$ is an infinite sequence of states
\[
  \rho=\rho_0\rho_1\rho_2\cdots
\]
such that $\rho_0=q_0$ and
$(\rho_i,\alpha_i,\rho_{i+1})\in\Delta$ for every $i\geq 0$.

The run is \textbf{accepting} if
\[
  \inf(\rho)\cap F\neq\emptyset.
\]
The word $\alpha$ is accepted if there exists at least one accepting run
on $\alpha$.
\end{definition}

For finite words, we ask where the run ends.  For infinite words, there
is no end.  B\"uchi acceptance instead asks whether the run visits some
accepting state infinitely many times.

This is also why ``membership'' changes character.  For a finite word we
can simulate the automaton to the last symbol and inspect the final
state.  For an $\omega$-word there is no last symbol to wait for.  The
decision procedures therefore reason about the finite graph of the
automaton, looking for a finite pattern that can repeat forever.

\begin{definition}[$\omega$-regular language]
\label{def:omega-regular-language}
A language $L\subseteq\Sigma^\omega$ is \textbf{$\omega$-regular} if it
is accepted by some B\"uchi automaton.
\end{definition}

This is the infinite-word counterpart of ``regular language''.  It tells
us that the class of properties expressible by B\"uchi automata has a
name and will turn out to have the same robust characterisations
(expressions, closure, decidable emptiness) that regular languages enjoy.

\begin{example}[Infinitely many acknowledgements]
\label{ex:buchi-infinitely-many}
Over $\Sigma=\{\texttt{req},\texttt{ack}\}$, the property
``\texttt{ack} occurs infinitely often'' is B\"uchi-recognizable.
Use two states: a non-accepting scanning state and an accepting state
entered whenever the current symbol is \texttt{ack}.  The accepting
state is visited infinitely often exactly when infinitely many
\texttt{ack} symbols appear.
\end{example}

\Cref{fig:infinitely-many-ack} shows the whole construction.  The
accepting state is not a final destination; it is a checkpoint that the
run must keep returning to.

\begin{figure}[ht]
\centering
\begin{tikzpicture}[
  node distance=3.2cm,
  state/.style={circle, draw, thick, minimum size=11mm, fill=blue!10},
  accepting/.style={circle, draw, double, thick, minimum size=11mm, fill=teal!10},
  arr/.style={-Stealth, thick},
]
  \node[state] (s) {$q$};
  \node[accepting, right=of s] (a) {$f$};
  \draw[arr] ([xshift=-1.0cm]s.west) -- (s);
  \draw[arr] (s) edge[loop above] node {\texttt{req}} (s);
  \draw[arr, bend left=18] (s) to node[above] {\texttt{ack}} (a);
  \draw[arr, bend left=18] (a) to node[below] {\texttt{req}} (s);
  \draw[arr] (a) edge[loop above] node {\texttt{ack}} (a);
\end{tikzpicture}
\caption{A B\"uchi automaton for infinitely many acknowledgements.}
\label{fig:infinitely-many-ack}
\end{figure}

\begin{example}[Parity between occurrences]
\label{ex:buchi-parity-thomas}
Here is a richer example, adapted from Wolfgang Thomas. Over the alphabet $\Sigma=\{a,b,c\}$, consider the language:
\begin{quote}
``Between any two consecutive occurrences of $a$, there is an even number of occurrences of $b$ and $c$ combined.''
\end{quote}
For instance, $a\,b\,c\,a$ is accepted (two symbols between the $a$'s), but $a\,b\,b\,b\,a$ is rejected (three symbols). Words with fewer than two occurrences of $a$ vacuously satisfy the condition.

The B\"uchi automaton uses its states to track the parity of the non-$a$ symbols since the last $a$:
\begin{itemize}[nosep]
  \item $q_0$ (initial, accepting): The empty prefix, or reading $b,c$ before the first $a$.
  \item $q_1$ (accepting): The last symbol was $a$, or we have seen an even number of $b/c$ symbols since the last $a$.
  \item $q_2$ (non-accepting): We have seen an odd number of $b/c$ symbols since the last $a$.
\end{itemize}
If the automaton reads an $a$ while in $q_2$, it transitions to a rejecting sink (omitted for clarity), because a forbidden odd-length block has been sealed by an $a$.
\end{example}

\Cref{fig:buchi-parity} makes the memory explicit: after the first $a$,
the automaton only needs to remember whether the current block length is
even or odd.

\begin{figure}[ht]
\centering
\begin{tikzpicture}[
  node distance=3.2cm,
  state/.style={circle, draw, thick, minimum size=11mm, fill=blue!10},
  accepting/.style={circle, draw, double, thick, minimum size=11mm, fill=teal!10},
  arr/.style={-Stealth, thick},
]
  \node[accepting] (q0) {$q_0$};
  \node[accepting, right=of q0] (q1) {$q_1$};
  \node[state, right=of q1] (q2) {$q_2$};
  \draw[arr] ([xshift=-1.0cm]q0.west) -- (q0);
  \draw[arr] (q0) edge[loop above] node {$b,c$} (q0);
  \draw[arr] (q0) -- node[above] {$a$} (q1);
  \draw[arr] (q1) edge[loop above] node {$a$} (q1);
  \draw[arr, bend left=18] (q1) to node[above] {$b,c$} (q2);
  \draw[arr, bend left=18] (q2) to node[below] {$b,c$} (q1);
\end{tikzpicture}
\caption{A B\"uchi automaton verifying the parity of $\{b,c\}$ blocks between occurrences of $a$.}
\label{fig:buchi-parity}
\end{figure}


% ============================================================
\section{Closure and Expressions}
\label{sec:omega-closure-expressions}
% ============================================================

The examples above show how to design individual automata.  For
verification, however, we also need to combine specifications.  If one
automaton describes ``requests are acknowledged infinitely often'' and
another describes a fairness assumption about the environment, then model
checking will ask about their conjunction, disjunction, or about a finite
initialisation phase followed by an infinite behaviour.  Closure
properties tell us that these combinations stay inside the same automata
model.

\begin{theorem}[Basic closure properties]
\label{thm:omega-basic-closure}
$\omega$-regular languages are closed under finite union and
intersection.  Moreover, if $U\subseteq\Sigma^*$ is regular and
$L\subseteq\Sigma^\omega$ is $\omega$-regular, then $U\cdot L$ is
$\omega$-regular; if $V\subseteq\Sigma^*$ is regular, then
$V^\omega$ is $\omega$-regular.
\end{theorem}

The constructions mirror the finite case, but one detail changes:
acceptance is about visiting final states infinitely often.  In product
constructions, we must ensure that both automata satisfy their B\"uchi
conditions, not just that both components are accepting at the same
single time.

\begin{formaldetails}[Construction ideas]
For $U\cdot L$, first run a finite-word automaton for $U$.  Whenever it
reaches an accepting state, nondeterministically switch to a B\"uchi
automaton for $L$.  The switch consumes no future symbols; it only marks
where the finite prefix ends.

For $V^\omega$, start from an automaton for the finite language $V$.
Assume, after the standard harmless normalization, that it has no
$\varepsilon$-moves and that no transition enters its initial state.
Whenever a transition would enter an accepting state after reading a
symbol $a$, duplicate that effect by sending the same $a$-transition
back to the initial state of the $V$-automaton, passing through a
B\"uchi accepting marker.  Each visit to the marker means that one
finite block from $V$ has just been completed.  Infinitely many marker
visits mean infinitely many completed $V$-blocks.

For intersection of two B\"uchi automata, the naive product with
$F_1\times F_2$ is too strict: one run may visit $F_1$ at even times and
the other may visit $F_2$ at odd times.  The correct product stores a
phase bit.  In phase $1$ it waits until the first component visits
$F_1$, then switches to phase $2$; in phase $2$ it waits until the
second component visits $F_2$, then switches back to phase $1$.  The
accepting states are the states that complete one of these obligations.
An accepting run must therefore complete both obligations infinitely
often.
\end{formaldetails}

\Cref{fig:buchi-intersection-phases} depicts only the control idea of the
intersection product.  The actual product also carries the two automaton
states; the phase records which acceptance obligation is currently being
checked.

\begin{figure}[ht]
\centering
\begin{tikzpicture}[
  phase/.style={draw, rounded corners=3pt, minimum width=25mm, minimum height=10mm, align=center},
  arr/.style={-Stealth, thick},
]
  \node[phase, fill=orange!10] (p1) {phase 1\\wait for $F_1$};
  \node[phase, fill=teal!10, right=32mm of p1] (p2) {phase 2\\wait for $F_2$};
  \draw[arr, bend left=18] (p1) to node[above] {hit $F_1$} (p2);
  \draw[arr, bend left=18] (p2) to node[below] {hit $F_2$} (p1);
  \node[below=10mm of p1, xshift=28mm, font=\small, align=center]
    {The product accepts only if both phases are completed infinitely often.};
\end{tikzpicture}
\caption{Asynchronous B\"uchi intersection is handled by alternating obligations.}
\label{fig:buchi-intersection-phases}
\end{figure}

\begin{definition}[$\omega$-regular expression]
\label{def:omega-regular-expression}
An \textbf{$\omega$-regular expression} denotes a finite union of
languages of the form
\[
  U_i V_i^\omega,
\]
where each $U_i,V_i\subseteq\Sigma^*$ is regular.
\end{definition}

Each term $U_i V_i^\omega$ describes a two-phase behaviour: a finite
initialisation from $U_i$, followed by an infinite repetition of blocks
from $V_i$.  The full expression is a finite union of such terms, so it
says ``the system does one of finitely many things: starts with some
prefix, then loops forever in some pattern''.  This is the infinite-word
analogue of a regular expression.

\begin{theorem}[B\"uchi-expression equivalence]
\label{thm:buchi-expression-equivalence}
A language $L\subseteq\Sigma^\omega$ is accepted by a B\"uchi automaton
if and only if it can be denoted by an $\omega$-regular expression.
\end{theorem}

\begin{proof}
We again prove both directions constructively.

\medskip
\noindent\textbf{Expression to automaton.}
It is enough to handle one term $UV^\omega$, because B\"uchi automata
are closed under finite union.  Let $A_U$ be an automaton for the regular
finite-word language $U$, and let $A_V$ be an automaton for $V$.  The
B\"uchi automaton first simulates $A_U$.  Whenever $A_U$ reaches an
accepting state, the automaton may non-deterministically switch to a
copy of $A_V$.  From then on, each time the $A_V$ copy accepts one block
from $V$, the automaton returns to the beginning of the $A_V$ copy and
passes through a B\"uchi accepting marker state.

Thus a run is accepting exactly when it reads one prefix from $U$ and
then completes infinitely many consecutive blocks from $V$.

\medskip
\noindent\textbf{Automaton to expression.}
Let $A=(Q,\Sigma,\Delta,q_0,F)$ be a B\"uchi automaton.  For states
$p,q\in Q$, let $W_{p,q}$ be the regular language of finite words that
label some path from $p$ to $q$ in the underlying finite automaton.  This
language is regular because it is accepted by the same finite graph with
initial state $p$ and accepting state $q$.

If $\alpha$ is accepted, then some accepting run visits at least one
state $f\in F$ infinitely many times.  Split the run at successive visits
to that same $f$.  The finite prefix before the first selected visit is
labelled by a word in $W_{q_0,f}$, and every later block is labelled by a
non-empty word in $W_{f,f}$.  Therefore
\[
  \alpha\in W_{q_0,f}(W_{f,f}\setminus\{\varepsilon\})^\omega.
\]
Taking the union over all $f\in F$ gives an $\omega$-regular expression
containing every accepted word.

Conversely, any word in one of these products describes a path from
$q_0$ to $f$, followed by infinitely many non-empty paths from $f$ back
to $f$.  Following those paths gives a run that visits $f$ infinitely
often, hence is accepting.  So the expression denotes exactly $L(A)$.
\end{proof}


% ============================================================
\section{Emptiness and Lassos}
\label{sec:buchi-emptiness}
% ============================================================

The most important algorithmic fact is that B\"uchi emptiness is still a
graph problem.

\begin{proposition}[B\"uchi emptiness]
\label{prop:buchi-emptiness}
Let $A$ be a B\"uchi automaton.  Then $L(A)\neq\emptyset$ if and only if
some accepting state is reachable from $q_0$ and belongs to a directed
cycle.
\end{proposition}

\begin{proof}
If an accepting state $f$ is reachable and lies on a cycle, follow the
path to $f$ once and then repeat the cycle forever.  This produces an
infinite run that visits $f$ infinitely often.

Conversely, any accepting run visits some accepting state $f$ infinitely
many times.  Between two visits to $f$ there is a non-empty finite path
from $f$ back to $f$, hence a cycle reachable from $q_0$.
\end{proof}

\aside{A counterexample in model checking
often has this shape: a finite prefix
followed by a repeated cycle.}

\Cref{fig:buchi-lasso} is the finite witness hidden inside the infinite
run: a prefix reaches an accepting cycle, and the cycle can then be
repeated without end.

\begin{figure}[ht]
\centering
\begin{tikzpicture}[
  node distance=2.2cm,
  state/.style={circle, draw, thick, minimum size=9mm, fill=blue!8},
  acc/.style={circle, draw, double, thick, minimum size=9mm, fill=teal!12},
  arr/.style={-Stealth, thick},
]
  \node[state] (q0) {$q_0$};
  \node[state, right=of q0] (q1) {$q_1$};
  \node[acc, right=of q1] (f) {$f$};
  \node[state, below=1.3cm of f] (q2) {$q_2$};
  \draw[arr] (q0) -- node[above] {$u$} (q1);
  \draw[arr] (q1) -- (f);
  \draw[arr] (f) -- node[right] {$v$} (q2);
  \draw[arr] (q2) -- (f);
  \node[below=7mm of q2, font=\small, align=center]
    {A reachable accepting cycle yields an infinite accepting run.};
\end{tikzpicture}
\caption{The lasso witness for non-emptiness: a finite prefix followed by a repeated cycle.}
\label{fig:buchi-lasso}
\end{figure}

\begin{definition}[Ultimately periodic word]
\label{def:ultimately-periodic}
An $\omega$-word is \textbf{ultimately periodic} if it has the form
\[
  uv^\omega
\]
for finite words $u\in\Sigma^*$ and non-empty $v\in\Sigma^*$.
\end{definition}

Every non-empty $\omega$-regular language contains an ultimately periodic
word.  The graph-theoretic lasso in \cref{prop:buchi-emptiness} is the
automata reason: a finite path gives $u$, and the repeated cycle gives
$v^\omega$.

\begin{example}[Languages without periodic words]
\label{ex:non-omega-regular-languages}
Consider the infinite word
\[
  \beta=a\,b\,a\,b^2\,a\,b^3\,a\,b^4\cdots .
\]
Both $a$ and $b$ occur infinitely often, but the gaps between consecutive
$a$'s keep growing. The word is not ultimately periodic: in an ultimately
periodic word, the distances between repeated occurrences are eventually
governed by one fixed period.

Can we build a non-empty language that is not $\omega$-regular?

We could simply take the singleton language $\{\beta\}$.  If it were
$\omega$-regular, by the lasso property, it would have to contain an
ultimately periodic word.  But its only word is $\beta$, which is not
ultimately periodic.  So $\{\beta\}$ is not $\omega$-regular.

That argument is correct, but it may feel too easy: a finite automaton for
one fixed aperiodic word already sounds impossible.  A more robust
counterexample should still have many words.

To make the obstruction robust, consider Wolfgang Thomas's example:
\[
  L_\beta = \{ u\beta : u\in\Sigma^* \}.
\]
This language is infinite: we can prepend any finite prefix to $\beta$.
Is $L_\beta$ $\omega$-regular?  If it were, it would have to contain an
ultimately periodic word.  But every word in $L_\beta$ shares the suffix
$\beta$, meaning it eventually behaves exactly like $\beta$ and cannot be
ultimately periodic.

Thus, $L_\beta$ is not $\omega$-regular.  The absence of an ultimately
periodic word is a structural barrier to $\omega$-regularity, even for
large infinite languages.
\end{example}

\subsection*{Notation Recap}
\addcontentsline{toc}{subsection}{Notation Recap}

\begin{center}
\begin{tabular}{@{}l l@{}}
\toprule
\textbf{Symbol} & \textbf{Meaning} \\
\midrule
$\Sigma^\omega$       & Set of all infinite words over $\Sigma$ \\
$\alpha[i,j]$         & Finite segment $\alpha_i\cdots\alpha_{j-1}$ \\
$\alpha[i,\omega]$    & Infinite suffix from position $i$ \\
$\inf(\alpha)$        & Symbols occurring infinitely often in $\alpha$ \\
$W^\omega$            & Infinite concatenation of blocks from $W\subseteq\Sigma^*$ \\
$\vec{W}$             & Vectorial closure: $\omega$-words whose prefixes hit $W$ infinitely often \\
$uv^\omega$           & Ultimately periodic word: finite prefix $u$, repeated cycle $v$ \\
$(Q,\Sigma,\Delta,q_0,F)$ & B\"uchi automaton tuple \\
$F$                   & Accepting states (must be visited infinitely often) \\
\bottomrule
\end{tabular}
\end{center}

\begin{summarybox}
\begin{itemize}[nosep]
  \item Infinite behaviours are modelled as words in $\Sigma^\omega$.
  \item B\"uchi automata accept runs that visit final states infinitely often.
  \item $\omega$-regular languages are exactly finite unions of $UV^\omega$.
  \item Emptiness reduces to finding a reachable accepting cycle.
  \item The lasso shape is the finite witness behind many infinite counterexamples.
\end{itemize}
\end{summarybox}

\section*{Exercises}
\addcontentsline{toc}{section}{Exercises}

\begin{chapex}[Computing an infinity set]
Let
\[
  \alpha = a\,b\,b\,a\,b\,b\,b\,a\,b\,b\,b\,b\cdots,
\]
where the number of $b$'s between consecutive $a$'s increases by one.
Compute $\inf(\alpha)$.  Is $\alpha$ ultimately periodic?
\end{chapex}

\begin{chapex}[B\"uchi acceptance]
Modify \cref{fig:infinitely-many-ack} to accept the property
``only finitely many \texttt{req} symbols occur''.  Hint: the automaton
will need a non-deterministic guess.
\end{chapex}

\begin{chapex}[Finding a lasso]
Given a B\"uchi automaton graph, explain how you would use strongly
connected components to decide emptiness.  Which SCCs are relevant?
\end{chapex}

\begin{chapex}[Writing an $\omega$-regular expression]
Over $\Sigma=\{\texttt{req},\texttt{ack}\}$, write an $\omega$-regular
expression for the language ``\texttt{ack} occurs infinitely often''.
Write a second expression for ``\texttt{ack} occurs only finitely
often''.
\end{chapex}
